% theme: gas_generation_properties
\MajorQuestion{気体の発生と性質}
\SetRowSpaceQanda

\Instruction{気体の集め方と空気の成分に関する次の問いに答えなさい。}
\QQ{水に溶けにくい気体を、水と置き換えて集める方法を\NamedBlank{water_displacement}という。}{}
\QQ{水に溶けやすく、空気より密度が大きい気体を集める方法を\NamedBlank{downward_displacement}という。}{}
\QQ{水に溶けやすく、空気より密度が小さい気体を集める方法を\NamedBlank{upward_displacement}という。}{}
\QQ{空気中に体積の割合で約78％含まれる気体を\NamedBlank{nitrogen}という。}{}
\QQ{空気中に体積の割合で約21％含まれる気体を\NamedBlank{oxygen}という。}{}

\Instruction{酸素の性質とつくり方に関する次の問いに答えなさい。}
\QQ{\RefBlank{oxygen}がもつ、ものを燃やすはたらきを何というか。}{}
\QQ{過酸化水素水に加えて、気体の発生を助ける黒色の物質は何か。}{}
\QQ{過酸化水素水の市販名を何というか。}{}
\QQ{過マンガン酸カリウムを加熱すると発生する気体は何か。}{}
\QQ{\RefBlank{oxygen}は水に溶けにくいため、どの方法で集めるか。}{}

\Instruction{二酸化炭素などの気体に関する次の問いに答えなさい。}
\QQ{石灰石にうすい塩酸を加えると発生する気体を\NamedBlank{carbon_dioxide}という。}{}
\QQ{\RefBlank{carbon_dioxide}を通すと、石灰水はどのように変化するか。}{}
\QQ{\RefBlank{carbon_dioxide}が水に少し溶けた水溶液は何性を示すか。}{}
\QQ{\RefBlank{carbon_dioxide}は、空気より密度が大きいか小さいか。}{}
\QQ{燃料の不完全燃焼で生じる、無色・無臭で有毒な気体は何か。}{}

\Instruction{水素とアンモニアの性質とつくり方に関する次の問いに答えなさい。}
\QQ{亜鉛にうすい塩酸を加えると発生する気体を\NamedBlank{hydrogen}という。}{}
\QQ{\RefBlank{hydrogen}が燃えるとできる物質は何か。}{}
\QQ{塩化アンモニウムと水酸化カルシウムを混ぜて熱すると発生する気体を\NamedBlank{ammonia}という。}{}
\QQ{\RefBlank{ammonia}は、水に溶けやすいか溶けにくいか。}{}
\QQ{\RefBlank{ammonia}の水溶液は何性を示すか。}{}
